\documentclass[11pt,a4paper]{article}
\usepackage[utf8]{inputenc}
\usepackage[T1]{fontenc}
\usepackage[ngerman]{babel}
\shorthandoff{"}
\usepackage{lmodern}
\usepackage{booktabs,tabularx,array}
\usepackage[table]{xcolor}
\usepackage[most]{tcolorbox}
\usepackage{microtype}
\usepackage{hyperref}
\usepackage{listings}
\usepackage[a4paper,left=18mm,right=18mm,top=15mm,bottom=15mm,headheight=0pt]{geometry}
\definecolor{navy}{HTML}{174A63}
\definecolor{blue}{HTML}{2E7698}
\definecolor{lightblue}{HTML}{EAF3F7}
\definecolor{lightgray}{HTML}{F2F3F4}
\definecolor{midgray}{HTML}{66727A}
\hypersetup{colorlinks=true,linkcolor=navy,urlcolor=blue,pdftitle={Sicherungsblatt Aufgabe 6.0 -- Programmierung wiederholen}}
\setlength{\parindent}{0pt}
\setlength{\parskip}{5pt}
\renewcommand{\familydefault}{\sfdefault}
\pagestyle{empty}
\newtcolorbox{definition}[1][]{enhanced,colback=lightblue,colframe=blue,
  boxrule=0.8pt,arc=1.5mm,left=3mm,right=3mm,top=1.5mm,bottom=1.5mm,
  fonttitle=\bfseries,title=Definition,#1}
\newtcolorbox{classcard}{enhanced,colback=white,colframe=navy,
  colbacktitle=lightblue,coltitle=navy,boxrule=0.8pt,arc=0mm,
  left=4mm,right=4mm,top=2mm,bottom=2mm,
  fonttitle=\bfseries,halign title=center,title=Hund,
  segmentation style={solid,navy}}
\newtcolorbox{objectcard}[1]{enhanced,colback=white,colframe=blue,
  colbacktitle=lightblue,coltitle=navy,boxrule=0.8pt,arc=2mm,
  left=4mm,right=4mm,top=2mm,bottom=2mm,
  fonttitle=\bfseries,halign title=center,title=\underline{\texttt{#1: Hund}}}
\newcommand{\sectionline}[1]{%
  \vspace{0.5mm}{\Large\bfseries\color{navy}#1}\par\nobreak\vspace{1mm}
  {\color{blue}\hrule height 0.8pt}\vspace{1mm}}
\newcommand{\sheettitle}[2]{%
  \begin{tcolorbox}[colback=navy,colframe=navy,arc=2mm,boxrule=0pt,left=5mm,right=5mm,top=2.5mm,bottom=2.5mm]
    {\color{white}\fontsize{19}{22}\selectfont\bfseries #1}\par
    \vspace{0.5mm}{\color{white}\large #2}
  \end{tcolorbox}}
\lstdefinestyle{java-sheet}{basicstyle=\ttfamily\small,
  columns=fullflexible,keepspaces=true,showstringspaces=false,
  aboveskip=2pt,belowskip=2pt}
\lstset{style=java-sheet}
\newcommand{\codecaption}[1]{{\small\bfseries\color{navy}#1}\par}

\begin{document}
\shorthandoff{"}
\sheettitle{Programmierung wiederholen}{Sicherungsblatt zu Aufgabe 6.0 \(\cdot\) Informatik 11 \(\cdot\) KI}

\sectionline{1. Klasse und Objekte}
\begin{classcard}
\texttt{private int alter}\\
\texttt{private String name}
\tcblower
\texttt{public void zeigeDaten()}\\
\texttt{public void belle()}\\
\texttt{public int gibAlter()}\\
\texttt{public void setzeAlter(int neuesAlter)}\\
\texttt{public void setzeDaten(int neuesAlter, String neuerName)}
\end{classcard}
\begin{center}
\textbf{Aufruf:} \texttt{petersHund.setzeAlter(6);}
\end{center}
\noindent\begin{minipage}{0.48\linewidth}\centering\textbf{Vor dem Aufruf}\end{minipage}\hfill
\begin{minipage}{0.48\linewidth}\centering\textbf{Nach dem Aufruf}\end{minipage}
\begin{tcbraster}[raster columns=2,raster equal height=rows,raster column skip=6mm,raster row skip=3mm]
\begin{objectcard}{petersHund}
\texttt{alter = 5}\\
\texttt{name = "Wuffti"}
\end{objectcard}
\begin{objectcard}{petersHund}
\texttt{alter = 6}\\
\texttt{name = "Wuffti"}
\end{objectcard}
\begin{objectcard}{inasHund}
\texttt{alter = 8}\\
\texttt{name = "Schnuffi"}
\end{objectcard}
\begin{objectcard}{inasHund}
\texttt{alter = 8}\\
\texttt{name = "Schnuffi"}
\end{objectcard}
\end{tcbraster}
Die Klasse legt Attribute und Methoden fest. Jedes Objekt besitzt eigene Attributwerte;
der Aufruf verändert nur Wuffti.

\sectionline{2. Methoden lesen und aufrufen}
\begin{definition}[title=Allgemeiner Methodenkopf]
\begin{lstlisting}
public Rueckgabetyp methodenname(Typ parametername)
\end{lstlisting}
\end{definition}
\begin{definition}[title=Beispiel aus Hund.java]
\begin{lstlisting}
public void setzeAlter(int neuesAlter)
\end{lstlisting}
{\small \texttt{void}: kein Rückgabewert \(\cdot\) \texttt{int}: Parametertyp \(\cdot\) \texttt{neuesAlter}: Parametername.}
\end{definition}
\begin{definition}[title=Punktnotation: Aufruf an einem bestimmten Objekt]
\begin{lstlisting}
objektvariable.methodenname(argumente);
petersHund.setzeAlter(6);
\end{lstlisting}
{\small \texttt{6} ist das \textbf{Argument}; sein Wert wird an \texttt{neuesAlter} übergeben.
Typ, Anzahl und Reihenfolge müssen passen. Ohne Argumente bleiben die Klammern leer.}
\end{definition}

\clearpage
\sheettitle{Programmierung wiederholen}{Sicherungsblatt zu Aufgabe 6.0 \(\cdot\) Informatik 11 \(\cdot\) KI}
\sectionline{3. Getter und Setter}
\begin{minipage}[t]{0.49\linewidth}
\vspace{0pt}
\begin{definition}[title=Getter: Attributwert zurückgeben]
\codecaption{Allgemeine Form für ein Attribut}
\begin{lstlisting}
public Typ gibAttribut() {
    return attribut;
}
\end{lstlisting}
\codecaption{Beispiel aus Hund.java}
\begin{lstlisting}
public int gibAlter() {
    return alter;
}
\end{lstlisting}
{\small \texttt{int}: Rückgabetyp \(\cdot\) \texttt{()}: kein Parameter. \texttt{return} gibt den Wert zurück.}
\end{definition}
\end{minipage}\hfill
\begin{minipage}[t]{0.49\linewidth}
\vspace{0pt}
\begin{definition}[title=Setter: Attributwert verändern]
\codecaption{Allgemeine Form für ein Attribut}
\begin{lstlisting}
public void setzeAttribut(Typ wert) {
    attribut = wert;
}
\end{lstlisting}
\codecaption{Beispiel aus Hund.java}
\begin{lstlisting}
public void setzeAlter(int neuesAlter) {
    alter = neuesAlter;
}
\end{lstlisting}
{\small Der übergebene Wert wird gespeichert. \texttt{void} bedeutet: kein Rückgabewert.}
\end{definition}
\end{minipage}
\begin{lstlisting}
petersHund.setzeAlter(7);
println(petersHund.gibAlter());  // 7
println(inasHund.gibAlter());    // 8
\end{lstlisting}
{\small Ein Getter gibt einen Wert zurück; erst \texttt{println(...)} zeigt ihn an.
Die Attribute sind \texttt{private}: Im Hauptprogramm erfolgt der Zugriff über zugängliche Methoden.}

\sectionline{4. Zwei Parameter und ein Aufruf ohne Argumente}
\begin{lstlisting}
public void setzeDaten(int neuesAlter, String neuerName)
inasHund.setzeDaten(9, "Flocke");
inasHund.zeigeDaten();
\end{lstlisting}
{\small Erst Alter (\texttt{int}), dann Name (\texttt{String}). Danach zeigt \texttt{zeigeDaten()} Flocke mit Alter 9.}
\end{document}
